\chapter{Элементарный объём и материя}
\label{ch:matter}

\section{Бинарная занятость и планкон}
\label{sec:planckon-floor0}

На каждой ячейке~$\PlanckCell$ вводится бинарная занятость $b(x)\in\{0,1\}$:
\textbf{планковская дырка} ($b=0$) или \textbf{планкон} ($b=1$).
Плотность материи принимает значения $0$ или планковская~$\rho_P$.
Первичный дефект --- одна ячейка.

\paragraph{Термин $\alpha_{\mathrm{fs}}$.}
\textbf{Постоянная тонкой структуры} на~$\Mlayer$ --- безразмерное EM-сопряжение заряда с вакуумом
(\cref{def:alpha-meaning}), число из дна носителя, не из $e,\planckHbar,c$.
Историческое имя из спектроскопии и lab-формула \eqref{eq:alpha-fs} --- T-export
(\cref{rem:alpha-fs-bridge,sec:alpha-name}).
Не путать с фазовым остатком $1/(4\pi)$ в насыщающей фазе (\cref{lem:vacuum-phase-residue}).

\section{Несжимаемая решёточная жидкость}

\begin{definition}[Модуль упругости $K_P$]
\label{def:kp}
$K_P=u_P$ --- планковское давление жидкости;
вращатель~$\Phi$ насыщается при $\rho\to u_P$.
\end{definition}

Отклонение плотности от фона вызывает упругий фазовый отклик;
избыток энергии уходит волной ($n=0$) или вихрем.
Уравнения Эйлера и Навье--Стокса на~$T$ --- макроосреднение.

\section{Истинный переносчик массы: квант \texorpdfstring{$\mathrm{Arg}$}{Arg}-натяжения и зигзаг (не поле Хиггса)}
\label{sec:true-mass-carrier}

В континуальной Стандартной модели массу частиц приписывают отдельному скалярному полю Хиггса:
предполагается, что исходно безмассовые поля «вязнут'' в непрерывном конденсате~$v$,
получая массу покоя через подгоняемые константы Юкавы.
На дискретном носителе~$\Mlayer$ эта картина перевёрнута с ног на голову:
\begin{enumerate}[leftmargin=*]
\item \textbf{На~$\Mlayer$ нет отдельного поля Хиггса и нет непрерывной «вязкости''.}
  Плотность материи в ячейке квантована жёстко: $\rho_{\mathrm{matter}}\in\{0,\rho_P\}$ ($b\in\{0,1\}$).
  Любой перенос по каузальным связям~$N_{12}$ всегда идёт с одной и той же тактовой скоростью
  $\caSignalSpeed=\caStepSpace/\htick=\sqrt{2}\,c$.
\item \textbf{Механизм инерции и массы покоя --- внутренний зигзаг в упругой ячейке~$K_P$.}
  Когда топологический дефект ($b=1$, вихрь) заперт в ячейке~$\hv$, каузальный перенос со скоростью~$\caSignalSpeed$
  натыкается на модуль упругости среды~$K_P$ и насыщающую фазу~$\Phi$.
  Фазовый поток
  \[
    \Delta\varphi_N = \mathrm{Arg}\!\left(\frac{\langle z\rangle_N}{z}\right)\in(-\pi,\pi]
  \]
  испытывает многократные внутренние отражения (зигзаг) в $\varepsilon$-окрестности узла~$\hv$,
  запирая энергию на месте. Макроскопический прибор на~$\Tlayer$, осредняя этот зигзаг,
  считывает его как эффективное замедление и массу покоя~$m_{\mathrm{rest}}$.
\item \textbf{Истинный первичный квант действия, энергии и массы на ячейке.}
  Из гейзенберговского пола фазы $\Delta\phi_{\min}=\tfrac12$ рад за один такт~$\htick=t_P/\sqrt{2}$
  следуют фундаментальные кванты одного $\mathrm{Arg}$-события:
  \begin{align}
    \ladderAction &= \planckHbar\,\Delta\phi_{\min} = \frac{\planckHbar}{2} \approx 5{,}273\times 10^{-35}\ \mathrm{Дж\cdot с},\\
    \ladderEnergy &= \frac{\ladderAction}{\htick} = \frac{E_P}{\sqrt{2}} \approx 1{,}383\times 10^{9}\ \mathrm{Дж} \approx 8{,}63\times 10^{27}\ \mathrm{эВ},\\
    m_{\arg} &= \frac{\ladderEnergy}{c^2} = \frac{m_P}{\sqrt{2}} \approx 1{,}539\times 10^{-8}\ \mathrm{кг}.
  \end{align}
  Деление $m_{\arg}$ на объём ячейки Вороного $\mathrm{d}V=\ell_P^3/\sqrt{2}$ немедленно замыкает
  планковскую плотность ячейки без подгоночных коэффициентов:
  \[
    \rho_{\mathrm{cell}} = \frac{m_{\arg}}{\mathrm{d}V} = \frac{m_P/\sqrt{2}}{\ell_P^3/\sqrt{2}} = \frac{m_P}{\ell_P^3} = \mu_P.
  \]
\item \textbf{Что такое бозон Хиггса на~$\Tlayer$.}
  Резонанс $m_H\approx 125{,}31\ \mathrm{ГэВ}$ ($\sim 1{,}25\times 10^{11}\ \mathrm{эВ}$), наблюдаемый на коллайдере, ---
  не первопричина массы, а \emph{вторичная макроскопическая пена} на $\mathrm{Arg}$-океане после прохождения
  через все $N_{\mathrm{hier}}=\lfloor B_{\hv}\rfloor-1=8$ битов иерархии ослабления $\alpha_{\mathrm{fs}}^8$
  ($v=\alpha_{\mathrm{fs}}^8 E_P\sqrt{2\pi}$, $m_H^{(\mathrm{bare})}=v/2$, гл.~\ref{ch:si-sm}).
  Один первичный квант $\ladderEnergy$ на~$\Mlayer$ мощнее массы бозона Хиггса в
  \[
    \frac{\ladderEnergy}{m_H c^2} \approx 6{,}89\times 10^{16}\ \text{раз}.
  \]
\end{enumerate}

\begin{table}[htbp]
  \centering
  \caption{Истинный переносчик инерции на~$\Mlayer$ против макро-резонанса Хиггса на~$\Tlayer$.}
  \label{tab:arg-vs-higgs}
  \begin{tabular}{@{}llll@{}}
    \toprule
    величина & формула & значение (СИ / эВ) & физический статус \\
    \midrule
    квант действия $\ladderAction$ & $\planckHbar/2$ & $5{,}27\times 10^{-35}\ \mathrm{Дж\cdot с}$ & 1 такт $\Delta\phi_{\min}=1/2$ на $\hv$ \\
    квант энергии $\ladderEnergy$ & $\ladderAction/\htick = E_P/\sqrt{2}$ & $1{,}38\times 10^{9}\ \mathrm{Дж}\ (\sim 8{,}6\times 10^{27}\ \mathrm{эВ})$ & первичный шаг энергии на $\Mlayer$ \\
    квант массы $m_{\arg}$ & $\ladderEnergy/c^2 = m_P/\sqrt{2}$ & $1{,}54\times 10^{-8}\ \mathrm{кг}$ & инерционный эквивалент $\mathrm{Arg}$-тика \\
    скорость переноса $\caSignalSpeed$ & $\caStepSpace/\htick = \sqrt{2}\,c$ & $4{,}24\times 10^{8}\ \mathrm{м/с}$ & каузальный шаг по связям $N_{12}$ \\
    бозон Хиггса $m_H$ ($\Tlayer$) & $\sqrt{2\lambda}\,v \sim \alpha_{\mathrm{fs}}^8 E_P$ & $125{,}31\ \mathrm{ГэВ}\ (\sim 1{,}25\times 10^{11}\ \mathrm{эВ})$ & вторичная пена ($\sim 10^{16}\times$ слабее $\ladderEnergy$) \\
    \bottomrule
  \end{tabular}
\end{table}

\section{Теорема о механической лестнице}
\label{sec:theorem-51}

\begin{theorem}[Лестница Ландау на ячейке]
\label{thm:51}
После теорем о дискретности и законе~$g$ (\cref{thm:discreteness,thm:g-law}) локальные механические величины на~$\hv$ образуют
единственную лестницу $\ladderAction\to \ladderMomentum\to L_0\to \ladderForce\to \ladderEnergy$
без трёх независимых подгоняемых параметров.
\end{theorem}

\begin{lemma}[5.1a]
\label{lem:51a}
\[
  \ladderAction=\planckHbar/2.
\]
\end{lemma}

\begin{proof}
\pfrom{минимального фазового события на такт (A5, дискретный ротор)}
\[
  \Delta\phi_{\min}=\tfrac12.
\]
Ненулевой осциллирующий фон задаёт минимальное изменение аргумента на~$\hv$ за такт~$\htick$.
\pfollows
\[
  \ladderAction=\planckHbar\,\Delta\phi_{\min}=\frac{\planckHbar}{2}.
\]
\end{proof}

\begin{lemma}[5.1b]
\label{lem:51b}
Из $\ladderAction$ и пары $(\caStepSpace,\htick)$: $\ladderMomentum=\ladderAction/\caStepSpace$, $\ladderEnergy=\ladderAction/\htick$, $m_{\arg}=\ladderEnergy/c^2$,
и потому $\ladderMomentum=m_{\arg} \caSignalSpeed/2$.
\end{lemma}

\begin{proof}
\pfrom{определения лестницы на $(\caStepSpace,\htick)$}
квант действия, отнесённый к шагу, есть импульс; тот же квант, отнесённый к такту, --- энергия:
\begin{align}
  \ladderMomentum &= \frac{\ladderAction}{\caStepSpace} = \frac{\planckHbar}{2\caStepSpace}, \\
  \ladderEnergy &= \frac{\ladderAction}{\htick}.
\end{align}
\pusing{тактовой скорости $\caSignalSpeed=\caStepSpace/\htick$}
\[
  \ladderEnergy=\ladderMomentum \caSignalSpeed.
\]
Масса аргумента определяется через макроскопическую скорость света из SI-моста:
$c=\kgeom \caSignalSpeed$ при $\kgeom=1/\sqrt{2}$, откуда $c^2=\caSignalSpeed^2/2$.
\ptransform{определения $m_{\arg}=\ladderEnergy/c^2$ и подстановки $c^2$}
\[
  m_{\arg}=\frac{\ladderEnergy}{c^2}=\frac{2\ladderEnergy}{\caSignalSpeed^2}=\frac{2\ladderMomentum}{\caSignalSpeed}.
\]
\pfollows $\ladderMomentum=m_{\arg} \caSignalSpeed/2$; множитель $1/2$ есть $\kgeom^2$.
\end{proof}

\begin{lemma}[5.1c]
\label{lem:51c}
$L_0=\ladderMomentum\caStepSpace=\ladderAction$, $g_M=\caSignalSpeed/(2\htick)$, $\ladderForce=\ladderMomentum/\htick=m_{\arg} g_M$ ---
соотношение $F=mg$ на одной ячейке~$\hv$.
\end{lemma}

\begin{proof}
\pfrom{определения момента на шаге}
\[
  L_0=\ladderMomentum\caStepSpace=\ladderAction.
\]
\pfrom{леммы~\ref{lem:51b}}
$\ladderMomentum=m_{\arg} v$ с $v=\caSignalSpeed/2$.
Набрать $v$ за один такт --- ускорение
\[
  g_M=\frac{\caSignalSpeed/2}{\htick}=\frac{\caSignalSpeed}{2\htick}.
\]
Сила одного кванта импульса за такт:
\[
  \ladderForce=\frac{\ladderMomentum}{\htick}.
\]
\psubst{$\ladderMomentum=m_{\arg} \caSignalSpeed/2$}
\[
  m_{\arg} g_M=\frac{\ladderMomentum}{\htick}=\ladderForce,
\]
то есть $F=mg$ на ячейке~$\hv$.
\end{proof}

\begin{lemma}[5.1d]
\label{lem:51d}
Ударный член $\lfloor\mathcal{N}\rfloor$ на~$\mathbb{Z}$ $\Rightarrow$ $\Delta p\in \ladderMomentum\cdot\mathbb{Z}$;
интеграл $\int\Im(z^*\nabla z)$ --- только макроописание на~$T$.
\end{lemma}

\begin{proof}
\pfrom{теоремы~\ref{thm:g-law}}
$\ladderMomentum=\ladderAction/\caStepSpace$ на связи.
\pusing{CL-4 (\cref{sec:congruence})}
удар даёт целое число квантов $\Delta\phi_{\mathrm{disc}}$ на такт, откуда
\[
  \Delta p=\ladderMomentum\cdot n,\qquad n\in\mathbb{Z}.
\]

\smallskip
\pfrom{отсутствия континуума на~$M$}
интеграл $\int\Im(z^*\nabla z)\,\mathrm{d}^3x$ не определён (нет меры Лебега).
\pusing{осреднения $\mathcal{B}$ (\cref{def:soft-readout})}
тот же поток читается на~$T$ как маделунговый ток.
\end{proof}

\begin{lemma}[5.1e]
\label{lem:51e}
$\ladderEnergy=\ladderMomentum \caSignalSpeed=\ladderForce\caStepSpace=L_0/\htick$, $E=n_E \ladderEnergy$, $n_E\in\mathbb{Z}$.
\end{lemma}

\begin{proof}
\pusing{лемм~\cref{lem:51b,lem:51c} и $\caSignalSpeed=\caStepSpace/\htick$}
\begin{align*}
  \ladderEnergy&=\frac{\ladderAction}{\htick}=\frac{\ladderMomentum\caStepSpace}{\htick}=\ladderMomentum \caSignalSpeed,\\
  \ladderForce\caStepSpace&=\frac{\ladderMomentum}{\htick}\cdot\caStepSpace=\ladderMomentum \caSignalSpeed,\\
  \frac{L_0}{\htick}&=\frac{\ladderAction}{\htick}=\ladderEnergy.
\end{align*}
\pfollows
три записи энергии совпадают.
\pfrom{CL-4 (\cref{sec:congruence})}
$n_E=\lfloor|\Phi|/\Delta\phi_{\mathrm{disc}}\rfloor\in\mathbb{Z}$ и $E=n_E \ladderEnergy$.
\end{proof}

\begin{proof}[Доказательство теоремы]
\pfrom{леммы~\ref{lem:51a}}
\[
  \ladderAction=\frac{\planckHbar}{2}.
\]
\pfrom{леммы~\ref{lem:51b}}
деление на $\caStepSpace$ и $\htick$ даёт $\ladderMomentum$ и $\ladderEnergy$;
масса через $c=\kgeom \caSignalSpeed$ даёт $\ladderMomentum=m_{\arg} \caSignalSpeed/2$.

\pfrom{леммы~\ref{lem:51c}}
сила одного кванта импульса за такт совпадает с $m_{\arg} g_M$ на ячейке~$\hv$.

\pfrom{леммы~\ref{lem:51d}}
приращение импульса на решётке целочисленно ($\Delta p\in\ladderMomentum\cdot\mathbb{Z}$),
а интеграл $\int\Im(z^*\nabla z)$ на~$M$ не определён.

\pfrom{леммы~\ref{lem:51e}}
три записи энергии совпадают:
\[
  \ladderEnergy=\ladderMomentum \caSignalSpeed=\ladderForce\caStepSpace=\frac{L_0}{\htick}.
\]
\phence свободного параметра, подгоняемого отдельно для импульса, силы и энергии, не остаётся.
\end{proof}

\begin{corollary}[Закон Ньютона на ячейке]
$\ladderMomentum=m_{\arg}(\caSignalSpeed/2)$, $\ladderForce=m_{\arg} g_M$, $\Delta F=\Delta m\cdot g$ при $\Delta m\in m_{\arg}\mathbb{Z}$.
\end{corollary}

\section{Термодинамические потенциалы на \texorpdfstring{$T$}{T}}
\label{sec:thermo-potentials}

Классические потенциалы относятся к макроуровню~$T$: на~$M$ нет температуры прибора
и нет гиббсовского ансамбля по~$g$ (\cref{prop:determinism}).
Их можно задать однозначно, если зафиксировать абсолютную энтропию по Планку
(\cref{ch:axiom-rationale}) и связать макроскопические $U$, $P$, $V$ с лестницей
Теорема~\ref{thm:51} после осреднения с биномиальным ядром (\cref{def:soft-readout}).

\paragraph{Внутренняя энергия.}
Из леммы~\ref{lem:51e} следует энергия на~$M$ дискретна: $E=\sum_x n_E(x)\,\ladderEnergy$,
$n_E\in\mathbb{Z}_{\ge 0}$.
На~$T$ в объёме~$V$ прибор видит гладкую плотность возбуждений
$\mathscr{n}_E(X)=\mathcal{B}*n_E$ и определяет
\begin{equation}
\label{eq:internal-energy}
  U(V)=\int_V \mathscr{n}_E(X)\,\ladderEnergy\,\mathrm{d}^3X
  = \int_V u(X)\,\mathrm{d}^3X,
  \qquad u=\mathscr{n}_E \ladderEnergy.
\end{equation}
Вакуум ($n_E=0$ везде) даёт $U=0$; это согласуется с $S\to 0$ при $T\to 0$,
если основное состояние невырождено ($W=1$ в формулировке Планка).

\paragraph{Давление и модуль.}
Модуль~$K_P=u_P$ (\cref{def:kp}) --- планковское давление несжимаемой решёточной жидкости:
вращатель насыщается при $\rho_E\to K_P$.
В термодинамике объём~$V$ --- макрообъём прибора; давление~$P$ --- сопряжённая
к~$V$ величина при медленных деформациях фона.
В натуральных единицах ($\caSignalSpeed=\caStepSpace/\htick=1$) равновесный вакуум ($\rho_E=0$)
лежит у нижней границы насыщения; для малых отклонений $\delta\rho=\rho_E\ll K_P$
\begin{equation}
\label{eq:pressure-linear}
  P \;\approx\; K_P\,\frac{\delta\rho}{K_P} = \delta\rho,
  \qquad
  \left.\frac{\partial P}{\partial V}\right|_{T}
  \sim -\frac{K_P}{V}\quad\text{при фиксированной массе.}
\end{equation}
Полная нелинейная связь $P(\rho_E)$ задаётся той же насыщающей фазой,
из которого в \cref{sec:gl-free-energy} получен член $\tfrac{\beta}{2}|\Phi|^4$
($\beta\equiv\omega_{\mathrm{macro}}$);
$\mathcal{F}_{\mathrm{GL}}$ описывает \emph{поле}~$\Phi$, а $P$ и~$U$ --- скаляры
субъекта наблюдения на~$T$.

\paragraph{Энтропия.}
На~$M$ энтропия фазы --- мера смешения аргументов~$z$
в пределах одной ячейки~$\hv$ при фиксированном~$Q$.
После~$\mathcal{B}$ прибор не различает микрофазу внутри вокселя:
макроэнтропия~$S$ --- функция $T$ и макропараметров, согласованная с
\[
  S = k_{\mathrm{B}}\ln W
\]
(Планк), где $W$ --- число равновесных \emph{макро}состояний, неразличимых readout'ом.
Условие $S\xrightarrow[T\to 0]{}0$ задаёт нулевую точку шкалы и делает
абсолютными потенциалы ниже.

\paragraph{Функции Гельмгольца и Гиббса.}
При фиксированных $T$ и~$V$ (канонический ансамбль на~$T$) определяют
\begin{equation}
\label{eq:helmholtz}
  F = U - TS.
\end{equation}
При фиксированных $T$ и~$P$ вводят энтальпию $H=U+PV$ и
\begin{equation}
\label{eq:gibbs}
  G = H - TS = U + PV - TS = F + PV.
\end{equation}
Связи $S=-(\partial F/\partial T)_V$, $P=-(\partial F/\partial V)_T$ --- стандартные
тождества при гладкости $U(T,V)$; модель не подменяет их новыми аксиомами,
а фиксирует \emph{микроскопические} якоря $\ladderEnergy$ и~$K_P$, из которых $U$ и~$P$
следуют после~\eqref{eq:internal-energy} и насыщения фазы.
Полевая свободная энергия~\eqref{eq:gl-free-energy} --- другой уровень описания:
$\mathcal{F}_{\mathrm{GL}}[\Phi]$ для неоднородного $\Phi(X)$, тогда как $F$ в~\eqref{eq:helmholtz}
--- скаляр равновесного объёма~$V$.
Химический потенциал $\mu=(\partial F/\partial N)_T$.

\section{Симметрии и дискретные заряды}

\begin{proposition}[Теорема Нётер на решётке]
На~$M$ сохраняются:
\begin{itemize}
\item заряд при $U(1)$: $Q=\sum|z|^2$;
\item импульс: $\mathrm{div}\,\pi=0$ по модулю~$\ladderMomentum$;
\item момент при $SO(2)$ изотropного~$N$: $L_z\in L_0\cdot\mathbb{Z}$;
\item энергия: $E\in \ladderEnergy\cdot\mathbb{Z}$;
\item топологический заряд: $n\in\mathbb{Z}$.
\end{itemize}
\end{proposition}

\section{Внутренний спектр и зонная логика}
\label{sec:planckon-internal-spectrum}

Три этажа вверх от планковского пола --- \fvolrefs{ch:floors-preface,ch:floor0,ch:floor1,ch:floor2}.
Здесь остаётся общая механика: атомные орбитали на~$T$ --- readout составного узла;
планкон несёт ту же логику квантования на~$\Mlayer$.

\section{Геометрический квант связи \texorpdfstring{$\kappa_{\mathrm{link}}=1/|N|$}{kappa=1/|N|}}

\begin{proposition}
На \latGCK\ $|N|=12$ $\Rightarrow$ $\kappa_{\mathrm{link}}=1/12=\gamma=\nu_{\mathrm{CA}}/(\caSignalSpeed\caStepSpace)$.
На гексагональном срезе $|N|=6$ $\Rightarrow$ $1/6$.
\end{proposition}

Спектральные методы калибруют макроуровень~$\layerMacro$.

\section{Фононы и спектр}

Периодичность~$\carrierLattice$ $\Rightarrow$ первая зона Бриллюэна;
$p=n \ladderMomentum+\planckHbar k$, $k$ дискретен.
Фонон вакуума:, $v_a=2\caSignalSpeed$.
\[
\omega(k)\approx v_a|k|
\]
Фонон вещества --- второй кристалл над микроуровнем (гл.~\ref{ch:carrier}).
